\section{Reconstrucción explícita desde los retornos completos}
\label{ret:seccion}

\subsection{Dominio de estados y extensión bilateral}

Se considera una coordenatización del dominio superviviente generado por
APP--TRIT--TPK. Un estado de esta coordenatización conserva las hojas operatorias,
la orientación trítica, el residuo, el cociente, el acarreo, la ruta,
la frontera, el registro de memoria y el origen del calendario.
La actualización actúa sobre este estado completo.

Sean \(X_0\) un espacio compacto de Hausdorff no vacío y
\(F:X_0\to X_0\) una actualización continua y sobreyectiva.
Estas son las hipótesis del resultado en la coordenatización considerada.
La extensión bilateral de la actualización es
\begin{equation}
 \mathscr X_F=
 \bigl\{(x_k)_{k\in\ZZ}\in X_0^{\ZZ}:
                  F(x_k)=x_{k+1}\text{ para todo }k\in\ZZ\bigr\},
 \label{ret:extension}
\end{equation}
con la topología inducida por el producto.

\begin{lemma}[Existencia y transporte bilateral]
\label{ret:existencia}
El espacio \(\mathscr X_F\) es compacto, de Hausdorff y no vacío.
Cada proyección coordenada \(\mathscr X_F\to X_0\) es sobreyectiva.
El desplazamiento
\[
 \sigma((x_k)_{k\in\ZZ})=(x_{k+1})_{k\in\ZZ}
\]
es un homeomorfismo de \(\mathscr X_F\).
\end{lemma}
\begin{proof}
Cada condición \(F(x_k)=x_{k+1}\) define un cerrado del producto
compacto de Hausdorff \(X_0^{\ZZ}\).
Fijados un índice \(k_0\) y un estado \(z\in X_0\), cualquier colección
finita de estas condiciones, junto con \(x_{k_0}=z\), admite solución:
se completan hacia delante mediante \(F\) y hacia atrás mediante
preimágenes, que existen por sobreyectividad. Las condiciones poseen,
por tanto, la propiedad de intersección finita. La compacidad produce
una historia bilateral con \(x_{k_0}=z\).

Los desplazamientos \(\sigma\) y
\(\sigma^{-1}((x_k))=(x_{k-1})\) conservan las relaciones de
\eqref{ret:extension}; ambos son continuos por su acción coordenada.
\end{proof}

\subsection{Muestreo trítico y reconstrucción de todos los pasos}

Fijado el origen del calendario, los instantes tipados de retorno son
\[
 \mathcal T_+
 =\{k\in\ZZ:1+(k\bmod 9)\in\{1,4,7\}\}=3\ZZ.
\]
En cada instante se conserva el estado completo \(x_k\).
Definimos el espacio de registros de retorno por
\begin{equation}
 \mathscr Y_+
 =\bigl\{(y_j)_{j\in\ZZ}\in X_0^{\ZZ}:
                 F^3(y_j)=y_{j+1}\text{ para todo }j\in\ZZ\bigr\}.
 \label{ret:registros}
\end{equation}

\begin{theorem}[Recuperación desde los retornos completos]
\label{ret:homeomorfismo}
La aplicación
\[
 r=\operatorname{Ret}_+:\mathscr X_F\longrightarrow\mathscr Y_+,
 \qquad r(x)_j=x_{3j},
\]
es un homeomorfismo. Su inversa está dada explícitamente por
\begin{equation}
 \boxed{
  \bigl(r^{-1}(y)\bigr)_{3j+s}=F^s(y_j),
  \qquad j\in\ZZ,\quad s\in\{0,1,2\}.}
 \label{ret:inversa}
\end{equation}
\end{theorem}
\begin{proof}
Para \(x\in\mathscr X_F\),
\(F^3(x_{3j})=x_{3j+3}\), de modo que \(r(x)\in\mathscr Y_+\).
Recíprocamente, cada entero \(k\) posee una única expresión
\(k=3j+s\), con \(j\in\ZZ\) y \(0\le s<3\), también cuando
\(k<0\). La fórmula \eqref{ret:inversa} define, por tanto, todas las
coordenadas de una sucesión \(x\).

Si \(s=0\) o \(s=1\), se tiene
\(F(x_{3j+s})=F^{s+1}(y_j)=x_{3j+s+1}\).
En la frontera entre dos bloques,
\[
 F(x_{3j+2})=F^3(y_j)=y_{j+1}=x_{3j+3}.
\]
Así, \(x\in\mathscr X_F\).
La lectura de sus coordenadas de índice \(3j\) devuelve \(y_j\).
En sentido inverso, las relaciones de \eqref{ret:extension} dan
\(x_{3j+s}=F^s(x_{3j})\), por lo que la reconstrucción de \(r(x)\)
devuelve \(x\) íntegramente.

La aplicación \(r\) es continua por ser una selección de coordenadas.
Cada coordenada de su inversa es la composición de una proyección
coordenada con una de las aplicaciones continuas \(I,F,F^2\).
La inversa es, por consiguiente, continua.
\end{proof}

Una coordenada discreta adicional, por ejemplo el índice \(m\in\ZZ\)
de una coordenatización de realización, se conserva mediante
\[
 \widetilde r:\ZZ\times\mathscr X_F
 \longrightarrow\ZZ\times\mathscr Y_+,
 \qquad \widetilde r(m,x)=(m,r(x)).
\]
Su inversa es \((m,y)\mapsto(m,r^{-1}(y))\).
El mismo resultado vale para un origen de calendario fijo \(a\):
se sustituyen \(x_{3j}\) y \(x_{3j+s}\) por
\(x_{a+3j}\) y \(x_{a+3j+s}\).

\subsection{Conjugación del transporte y conservación de las lecturas}

Sea \(\sigma_+(y)_j=y_{j+1}\) el desplazamiento de los registros.
La reconstrucción anterior da
\begin{equation}
 r\sigma^3=\sigma_+r,\qquad
 (r\sigma r^{-1}(y))_j=F(y_j).
 \label{ret:conjugacion}
\end{equation}
La aplicación de la segunda igualdad es reversible en
\(\mathscr Y_+\); su inversa es
\[
 y\longmapsto\bigl(F^2(y_{j-1})\bigr)_{j\in\ZZ}.
\]
En efecto, las dos composiciones se reducen a
\(F^3(y_{j-1})=y_j\).
Así, la reversibilidad de la historia bilateral se conserva aun
cuando la actualización \(F\) sobre un estado aislado sea no inyectiva.
Si la vuelta nonádica se realiza como \(\Gamma_9=\sigma^9\), entonces
\begin{equation}
 r\Gamma_9=\sigma_+^3r.
 \label{ret:vuelta-nonadica}
\end{equation}

\begin{corollary}[Lecturas conjuntas y transporte recuperado]
\label{ret:lecturas}
Sean \(a:\mathscr X_F\to\mathcal R\) y
\(q:\mathscr X_F\to\mathcal E\) dos lectores continuos del mismo
estado, y sea \(d:\mathscr X_F\to\mathscr X_F\) un homeomorfismo del
dominio calibrado. Entonces
\[
 a_+=a\circ r^{-1},\qquad
 q_+=q\circ r^{-1},\qquad
 d_+=r\circ d\circ r^{-1}
\]
son continuos, \(d_+\) es un homeomorfismo y
\begin{equation}
 (a,q)=(a_+,q_+)\circ r,\qquad rd=d_+r.
 \label{ret:factorizacion}
\end{equation}
Para todo entero \(N>0\),
\[
 d^Nx=x\quad\Longleftrightarrow\quad d_+^Nr(x)=r(x).
\]
\end{corollary}
\begin{proof}
Se componen las aplicaciones con el homeomorfismo del
\cref{ret:homeomorfismo}. La inversa de \(d_+\) es
\(rd^{-1}r^{-1}\). La identidad
\(d_+^Nr=rd^N\), junto con la inyectividad de \(r\), prueba la
última equivalencia.
\end{proof}

La aplicación a la doble lectura toma
\(\mathcal R=R_{\rm HMT}\) y el codominio incidencial construido en
\cref{exc:doble-lectura,exc:limite-inverso}.
El resultado establece una presentación recuperable de las mismas
historias, de sus lectores y de su transporte. Su objeto es esta
equivalencia dinámica de presentaciones.
